\section{Relating LIBERO-Plus Responses to Commands and Task Outcomes}
\label{sec:supp-plus}

This section describes the $\pi_{0.5}$ LIBERO-Plus recordings to which the shared directions of main Sec.~\ref*{M-sec:shared} are fitted (first subsection); the rest is an additional study that the main paper does not rely on. Camera and Sensor Noise cost success; perturbed runs end farther from their clean runs than a clean rerun does, even when both succeed; and a prespecified test of whether response size identifies lost successes is inconclusive. This is consistent with main Sec.~\ref*{M-sec:sensitivity}, whose outcome figures use the separate paired executions (Sec.~\ref{sec:five-arm}); related analyses: Secs.~\ref{sec:current-basis-readout} and~\ref{sec:supp-bounded-internal}.

\subsection{Task outcomes of the recorded pairs}

\paragraph{The recordings: 536 pairs on 40 tasks.}
The recordings contain 268 settings, each under two policy-noise draws: 536 clean--perturbed pairs on 40 tasks in the LIBERO-Plus categories Background Textures, Camera Viewpoints, Light Conditions and Sensor Noise. Two inspected Light settings also change texture assets, so Light denotes the benchmark's prescribed settings, not an isolated lighting factor. Each pair has its own executed clean run; the 80 clean-run recipes fix only initialization and noise settings, not 80 trajectories shared across pairs. Outcomes (clean$\to$perturbed) are 447 joint successes (SS), 56 lost successes (SF), 20 gained successes (FS) and 13 joint failures (FF). Pairs are nested within tasks.

\paragraph{The common-task comparison: 404 pairs on 28 tasks.}
All quantities below use one population, so that they are directly comparable: the 28 tasks shared by all four families, without 18 Light pairs lacking a LIBERO-Plus difficulty level, leaving 202 settings and 404 pairs, those of main Fig.~\ref*{M-fig:token-structure} (Table~\ref{tab:plus-population}). The estimates describe the selected settings; full-panel summaries and full-benchmark rankings refer to different populations.

\begin{table}[!ht]
\caption{\textbf{Pairs per family in the full panel and on the common tasks.} $\pi_{0.5}$ LIBERO-Plus recordings: pairs/tasks in the full panel, pairs/settings on the common tasks, and common tasks per family. The full Light panel includes 18 pairs without a difficulty level. The common tasks are shared, but families still differ in settings, difficulty levels and number of settings per task.}

\centering\small
\setlength{\tabcolsep}{2.5pt}
\begin{tabular}{lrrr}
\toprule
Family & Full: pairs/tasks & Common: pairs/settings & Tasks\\
\midrule
Background & 118/34 & 96/48 & 28\\
Camera & 112/32 & 102/51 & 28\\
Light & 152/40 & 96/48 & 28\\
Sensor Noise & 154/39 & 110/55 & 28\\
\bottomrule
\end{tabular}

\label{tab:plus-population}
\end{table}

\paragraph{Camera and Sensor Noise cost success.}
On the common tasks, the success change (perturbed minus clean) is $+2.679$, $-17.857$, $+2.679$ and $-12.500$ percentage points for Background, Camera, Light and Sensor Noise, with pointwise 95\% intervals $[-5.357,11.607]$, $[-30.379,-6.250]$, $[-7.143,13.393]$ and $[-22.321,-4.464]$, and simultaneous 95\% intervals over the four families (Bonferroni-adjusted percentiles) of $[-7.143,13.393]$, $[-34.821,-3.571]$, $[-8.929,16.071]$ and $[-25.893,-2.679]$. The Camera and Sensor Noise intervals exclude zero; those of Background and Light, with positive point estimates, include both gains and losses. Each change is measured against the family's own clean runs and mix of difficulty levels. Estimates average the two noise draws within a setting, settings within a task--family, and then tasks equally; intervals use 10,000 bootstrap draws of whole tasks.

\subsection{Definitions of the response and attention measures}
\label{sec:plus-sites}
Absolute response size measures a difference, relative size compares it with the clean tensor, and attention total variation (TV) measures a change in how attention is allocated.

\paragraph{Recording sites.}
The first policy query is $q=0$; $j$ indexes generation steps and $p$ plan positions. We record L1, L9 and L18 (zero-based block indices 0, 8, and 17). Prefix measurements replay the prefix computation with masks for valid tokens and defined queries. $H$ is the block input and $Z$ the pre-attention RMSNorm output supplied to attention, so a larger $H$ response at depth need not mean a larger change in what attention receives. Action-expert (AE) hidden states are recorded during the policy's own generation and keep the ten ordered plan positions at each generation step. Channel indices belong to their own representation and layer; equal indices in different spaces do not imply a shared direction.

\paragraph{Absolute and relative responses can move in opposite directions.}
For a selected set of $N$ scalar coordinates, let $X_R$ and $X_T$ denote the clean and perturbed tensors and $\Delta X=X_T-X_R$. We report
\begin{align}
 d_{\rm RMS}(X)&=\|\Delta X\|_F/\sqrt{N},\\
 r(X)&=\|\Delta X\|_F/\|X_R\|_F.
\end{align}
Here the absolute AE L18 response falls from the first to the last generation step in every family--suite cell, while the relative response rises, also in every cell, because the clean norm falls. A relative trend therefore mixes response size with the clean scale. Ratios are formed per pair before averaging and are undefined for a zero clean norm; a ratio of family-average norms is a different estimator.

\paragraph{Attention change is measured over all visible keys.}
The attention probabilities $P$ are not stored during the policy's run but recomputed at the recorded generation states. For each valid action query and attention head, the total variation over all visible keys is
\begin{equation}
 \operatorname{TV}(P_T,P_R)=\tfrac12\sum_{k\in\mathcal V}|P_T(k)-P_R(k)|,
\end{equation}
where $\mathcal V$ is the common visible-key set. We average over valid queries and heads and keep the probability mass on visual, text and other visible keys. The regional net mass change sums the signed probability change in a region; the regional conditional-distribution change first normalizes each run by its regional mass. Neither substitutes for full-key TV. The text region contains text tokens; the recorded configuration does not append the numerical joint state as discrete text.

\paragraph{Attention output combines allocation and transformed values.}
The content readout recomputes the attention output from prefix and action-side values, after projection, before the residual gate. Values and their weighted combination can change even when probability shifts decrease. This recomputation measures each branch at its saved generation state, separately from the recording of the policy's own generation (Table~\ref{tab:recording-roles}).

\paragraph{Commands and motion keep their own units.}
The output record separates the 32-coordinate policy actions from the seven-coordinate LIBERO commands after the adapter. We score translation, rotation and gripper components over the first five and over all ten plan positions. Command RMS uses saved command units; end-effector (EEF) distance uses the simulator's length units, without conversion to meters or millimeters. Orientation distance is the quaternion geodesic angle in radians. Gripper-state distance is the Euclidean difference of two recorded joint coordinates, not an opening width.

\subsection{Motion compared with the clean rerun}
\label{sec:plus-motion}
After 60 controls, perturbed runs are farther from their clean runs than a clean rerun is in every family, and under Camera and Sensor Noise also when both runs succeed. Two references give the separation a scale: the clean rerun, which changes only the policy noise, and a repeat of the clean run with the same inputs and noise.

\paragraph{Command differences are followed by motion.}
All 404 common-task pairs start at zero EEF-position, orientation and gripper-state distance, and all issue their first five commands. After 5 controls, the task-equal EEF distances are 0.002911, 0.008362, 0.002027 and 0.004108 for Background, Camera, Light and Sensor Noise, with pointwise intervals above zero. After 60 controls the means are 0.024468, 0.041819, 0.018961 and 0.030180. Each family keeps all 28 tasks over this window; later differences also reflect the different states the runs reach.

\paragraph{After 60 controls, separation exceeds the clean rerun.}
Let $\overline d_{TR,f}(60)$ be the task-equal clean--perturbed EEF distance for family $f$ and $\overline d_{RR,f}(60)$ that between the two clean runs. After 60 controls their ratios (above 1: farther apart than a clean rerun) are 1.51, 2.39, 1.17 and 1.68 for Background, Camera, Light and Sensor Noise. Each ratio divides two family means, not an average of per-task or per-setting ratios, and has no estimated interval. The paired distance-difference intervals (2,000 bootstrap draws of the fixed 28 tasks, pointwise 95\%) are in the released analysis package.

\paragraph{Extra motion and net success measure different effects.}
Background has a positive interval for motion beyond the clean rerun, yet its success interval includes zero. With the repeat as reference, the Background/Camera/Light/Sensor Noise reference distances are 0.002518/0.001873/0.002512/0.002565 on 28/26/26/28 tasks; the two references keep their own populations and means.

\paragraph{Runs that both succeed still move apart.}
Outcome groups refer to the case's own clean--perturbed pair: ALL includes SS, SF, FS and FF, and SS and SF use that pair's final outcomes, not those of every repeat. After 60 controls, with the clean rerun as reference, Camera has $d_{TR}=0.04360$ in SS (75 cases, 25 tasks) and $0.04291$ in SF (23 cases, 12 tasks); Sensor Noise has $0.02872$ in SS (91 cases, 27 tasks) and $0.03456$ in SF (16 cases, 10 tasks). The SS intervals for motion beyond the clean rerun are positive for Camera, $0.02619\ [0.01678,0.03596]$, and Sensor Noise, $0.01047\ [0.00181,0.01859]$, while those of Background, $0.00730\ [-0.00019,0.01414]$, and Light, $0.00228\ [-0.00348,0.00756]$, include zero. These outcome-defined groups have different support. Both runs can thus succeed along different paths. Beyond 60 controls, finished episodes shrink the paired set, so later values need their task and case counts.

\paragraph{The two references.}
The repeat matches the recorded initial inputs, configuration, random-number state and policy noise; for each case it averages all eligible, independently executed repeats, whose remaining variation need not be pure floating-point error. The clean rerun compares the two policy-noise draws within each setting. All 202 settings support the clean-rerun comparison; the repeat supports 396 of the 404 cases. Both references keep the case's clean--perturbed pair.

\subsection{Does response size identify lost successes?}
\label{sec:plus-p0}
This prespecified test asks whether response size ranks lost successes (SF) above joint successes (SS) within task. Its adjusted intervals cross the thresholds of both prespecified criteria, for discrimination and for its absence, so it supports neither conclusion; the secondary pooled AUROCs are modest. This is consistent with main Sec.~\ref*{M-sec:sensitivity} but does not by itself establish it.

\paragraph{Four directed scores fixed before the test.}
The comparison starts from all 518 labelled pairs, not only the 28 shared tasks, and uses four fixed scores: the relative AE L18 response over the ten plan positions at the first generation step, the relative prefix main-camera L18 response, the L18 full-key attention TV at the first generation step, and the RMS change of the first five commands (translation). Larger values always rank as a more likely lost success; this direction was fixed in advance and kept after the results were seen.

\paragraph{Within-task discrimination with prespecified criteria.}
The primary AUROC compares SF with SS within each task, with half credit for ties, weighting eligible tasks equally; a task lacking either outcome has no primary AUROC. Camera contributes ten eligible tasks with 15 SF and 23 SS cases; Sensor Noise 13 with 18 SF and 32 SS cases. The positive criterion requires the same score's adjusted lower bound to exceed 0.60 under both Camera and Sensor Noise; the criterion of no practical discrimination requires all eight adjusted upper bounds to lie below 0.70. The 98.75\% intervals for this restricted support (20,000 whole-task bootstrap draws shared across quantities and families, Bonferroni-adjusted over the four scores within each family) are in the released analysis package.

\paragraph{A secondary pooled comparison uses different support.}
The pooled secondary analysis keeps all SF and SS cases, takes the logarithm of each score with a fixed $10^{-12}$ floor and subtracts its task--family mean, so that task and family differences do not drive the result, before computing a pooled AUROC. Camera has 25 SF and 83 SS cases, with AUROCs 0.623/0.636/0.587/0.648 for the four scores; Sensor Noise has 22 SF and 125 SS, with 0.548/0.527/0.581/0.608. The adjusted intervals of the command score are [0.501,0.817] and [0.408,0.802]. Comparing higher with lower LIBERO-Plus difficulty levels under Camera and Sensor Noise, several responses increase while within-task success-change intervals cross zero.
